\section{Introduction}

On a single smooth projective Calabi--Yau threefold, the scalar
local factors in the universal series of effective curves can occur
with exponents of arbitrarily large denominator. This already happens
in classes of divisibility three. To see what is being factored, assign
to each effective integral homology class the unit of its compact
space of cycles. Form the generating series of these units. A choice
of local model for the multiples of an embedded curve gives a factor
for each class and genus. Scaling the two normal directions with
opposite weights in a balanced bundle gives integral exponents
\cite{WueProduct,WueGenerators}. Scaling both directions with weight one gives
the scalar factors introduced by Pardon \cite{Pardon}.
The same spaces of cycles enter both products.

The compact example of \cite[Theorem~1.1]{WueCounter} concerns a
threefold containing a Hirzebruch surface. A primitive class on the surface has genus-three multiplicity
$-14$. At three times this class the genus-ten scalar exponent has
fractional part $2/3$. There are two
questions here. What denominator does the local class of a tripled
curve carry? Can a single compact target realize increasingly large
denominators once the covering degree has been fixed?

For the local question let $C$ be a smooth projective curve of genus
$g$ and $N$ a rank-two normal bundle of degree $2g-2$. The tripled zero
section is the effective cycle $3[C]$ in the total space of $N$.
An evaluation at a point of $C$ applies a fiber covector to the three
sheets of a cubic cycle and takes a homogeneous symmetric polynomial
in the three values. Its degree is its weight under scalar dilation.
A tuple with compact common zero locus gives an integral supported
Thom class. Dividing it by the product of the weights gives Pardon's
scalar class $x_{g,3}$ \cite[Definition~4.7]{Pardon}. Its denominator
is the least positive integer that takes it into $H_{\ZZ}$, the part
of the integral algebra of semi-Fano cycles in first Chern number and
virtual dimension zero. Here semi-Fano means that each irreducible
component has nonnegative first Chern number. We use the integral
generators of \cite[Theorem~1.1]{WueGenerators}, obtained by giving opposite
weights to two normal lines of degree $g-1$ whose product is $K_C$.

\begin{theorem}\label{thm:cubic}
In genera zero and one the scalar and opposite-weight classes coincide
in every degree. For $g\ge2$ the denominator of the scalar triple is
\[
 \operatorname{den}(x_{g,3})=
 \begin{cases}
  3^k, &g=2k+1,\\[2pt]
  3^{\max\{0,\,k-2-v_3(2k+1)\}}, &g=2k+2,
 \end{cases}
\]
where $v_3$ is the valuation at three. For a fixed curve of genus
$g\ge2$, the least product of weights, over all rank-two bundles of
degree $2g-2$ and all tuples of homogeneous single-fiber evaluations
with compact common zero locus on cubic cycles, is
$3^{\lfloor(g-1)/2\rfloor}$. A general stable bundle with determinant
$K_C$ attains the minimum.
\end{theorem}

For example, on a genus-four curve the best choice of normal bundle
allows a single evaluation of weight three. The resulting supported
class is divisible by three in $H_{\ZZ}$, so the normalized class is
integral. In genus five the minimum is nine and the class itself has
denominator nine. Computing a minimum weight product thus gives a
sharp geometric bound on the normalization, while computing the class
denominator requires the divisibility of the numerator as well.
For use below we write $s_{g-1}$ for the exponent of three in
$\operatorname{den}(x_{g,3})$ and put $s_0=0$.

Let $e_{\beta,h}(Y)$ denote the exponent of the scalar genus-$h$
factor at an effective class $\beta$. On our compact threefold we
find infinitely many of these individual exponents to which the
odd-genus formula applies exactly.

\begin{theorem}\label{thm:arith-compact}
There is a smooth projective Calabi--Yau hypersurface $Y$ containing
a first Hirzebruch surface with negative section $s$ and ruling $f$.
For each $a=12j+3$, $j\ge0$, put $\alpha_a=as+(a+1)f$ and
$k=(a(a-1)/2-1)/2$. The class $\alpha_a$ is primitive and
$e_{3\alpha_a,9k+1}(Y)$ has exact denominator $3^k$. The first of
these exponents, in genus ten, has fractional part $2/3$.
\end{theorem}

The local Gromov--Witten theory of Bryan and Pandharipande \cite{BP}
computes numerical residues for two normal weights. A numerical
residue alone does not identify a class in a universal cycle algebra.
To make the comparison, we vary a balanced sum of two normal lines
in a family of algebraic targets and impose conditions on the
characteristic polynomials of the sheets. At prime covering degree
the nonconstant coefficients in the ratio of the two weights can be
computed on embedded images, after deformation at each prescribed
bound on source genus. Intersection on the cycle space counts each
of these images with its sign. On the same family, stable maps with
that image include maps with contracted components, whose degree-one
series is $(2\sin(u/2))^{2h-2}$ for an image of genus $h$.
Theorem~\ref{thm:main} consequently identifies the difference of the
scalar and opposite-weight logarithms in prime degree. For triples,
the lower-degree terms are integral, so the cubic trace in
\S\ref{sec:arithmetic} determines the denominator of the full class.
A stable-pair calculation on the Hirzebruch surface singles out a
primitive genus with multiplicity prime to three. Its contribution
survives in the scalar exponent of Theorem~\ref{thm:arith-compact}.
The general formula in \S\ref{sec:compact} sums over all effective
integral roots of the triple class.

The genus-ten example distinguishes three measurements of the triple.
For a genus-three core the scalar-minus-opposite-weight cubic
logarithm has coefficient $4720/3$ at $(2\sin(u/2))^{18}$. Changing
variables and taking the coefficient of $u^{18}$ instead gives the
fixed source-genus-ten residue $23593/5$. It is the former coefficient
that multiplies $-14$ in the change of the compact exponents.
A different comparison uses the point class of an isolated triple in a
nonlinear neighborhood of the curve. For a general stable normal bundle of
canonical determinant there are eight one-parameter families of
embedded cyclic triples in genus three. A general quadratic term
obstructs all eight, and the isolated triple exceeds the scalar class
by $8x_{10,1}/3$, where $x_{10,1}$ is the class of a simple genus-ten
curve (Theorems~\ref{thm:classification} and~\ref{thm:isolated}).
This describes the moving covers behind the failure of triple
rigidity studied in \cite{BP06}.

In composite degree the characteristic polynomials can be proper
powers. We prove the comparison outside the weights where these
powers occur and compute the correction from doubled images. This
determines the scalar-to-opposite-weight conversion in degree four
for every core genus and in degree six for a genus-two core
(\S\ref{sec:composite}). Appendices~\ref{app:split}
and~\ref{app:pencils} treat individual split cubic weights and
the contribution of cuspidal cubics over a genus-five curve.

\subsection*{Acknowledgements}
I thank John Pardon for a helpful discussion of equivariant local curve classes.
